\subsection{Finite presentations and a decision obstruction}
Finite operation tables and finite presentations are different inputs.
The following reduction makes the ordinary ring difficulty explicit.
\begin{proposition}\label{prop:projective-decision-reduction}
A finite presentation of any unital associative ring $R$ can be
transformed effectively into a finite presentation of the inverse
semiring $T=R\times\B$, with a named idempotent $c=(1_R,0_\B)$.
Its additive-idempotent support is exactly $\B$. The cyclic projective
$Tc$ has singleton additive support and is zero exactly when $R=0$.
Thus uniform zero testing, and hence isomorphism testing, for specified
finite projectives over finitely presented inverse semirings is at
least as hard as triviality of finitely presented unital rings, even
with fixed Boolean support and a one-by-one idempotent matrix.
\end{proposition}
\begin{proof}
Write a finite ring presentation with generators $x_i$ and polynomial
relations $f_j(x;1)=0$. In the finitary inverse-semiring signature,
put $e=1+(-1)$ and use generators $c,z_i$ with relations
\[
 c+(-c)=0,\quad z_i+(-z_i)=0,\quad c+e=1,
 \quad cz_i=z_i=z_ic,\quad f_j(z;c)=0.
\]
The last relations replace every occurrence of the ring unit by $c$.
Distributivity and uniqueness of additive inverses give
$ce=ec=0$ and $ez_i=z_ie=0$. Multiplying $c+e=1$ by $c$ gives
$c^2=c$. Also $e+e=e$, $-e=e$ and $e^2=e$.
The bottom additive group therefore contains the ring generated by
$c,z_i$, with unit $c$, and receives the presented ring $R$.
Every term has the form $a+\epsilon e$, where $a$ is in that ring
and $\epsilon\in\{0,1\}$: these forms contain the generators and
$1=c+e$, and are closed under all operations because mixed products
with $e$ vanish.

The relations hold in $R\times\B$ under
$c=(1_R,0)$ and $z_i=(x_i,0)$. Conversely, the ring map into the
bottom group extends to
$(r,\epsilon)\mapsto r+\epsilon e$ in the presented inverse semiring.
The two maps are inverse on generators and normal forms, proving
that the finite presentation is exactly $R\times\B$, also when
$R=0$. Its two additive idempotents are $(0_R,0)$ and $(0_R,1)$.

Right multiplication by $c$ is an idempotent endomorphism of the
free left module $T$. Its image is $Tc=R\times\{0\}$, a cyclic
projective with a single additive idempotent. It vanishes exactly
when $R$ does. The construction supplies the certified idempotent
matrix $[c]$; no algorithm is asked to discover its idempotence.
Comparison with the zero projective also proves the isomorphism-test
reduction.
\end{proof}

The unit convention in an ordinary undecidability input must be
checked. Bokut's original presentation convention uses free
algebras without a multiplicative identity
\cite[Section~2, p.~137]{Bokut}. The following explicit bridge is
useful independently of that source application.
\begin{lemma}\label{lem:unital-idempotent-bridge}
Let $k$ be a finite field, $A$ a finitely presented possibly nonunital
$k$-algebra, and $e\in A$ a named idempotent. Let $S=k\oplus A$ be
its unitization and put
\[
 B=(S*_k k\langle u,v\rangle)/(uv-1,\ vu-e).
\]
The free product imposes no commutation of $u,v$ with $S$.
This is an effective finite unital presentation. Moreover, $B=0$
if and only if $e=0$ in $A$. If $e\ne0$, the natural map $S\to B$
is injective.
\end{lemma}
\begin{proof}
The same generators and nonconstant relations that present $A$,
now interpreted in the free unital algebra, present $S$.
Adjoining $u,v$ and the two displayed relations is effective.
If $e=0$, then $1=(uv)^2=u(vu)v=0$ in $B$.

Suppose $e\ne0$. Let $V$ be a countable direct sum of copies of the
left regular $S$-module, regarded as a $k$-vector space. Left
multiplication gives a faithful unital representation
$\lambda:S\to\operatorname{End}_k(V)$.
The algebra $S$ is spanned by countably many words. Both $V$ and
$eV$ have countably infinite dimension, since $S$ and $eS$ are
nonzero. Choose a vector-space isomorphism $w:V\to eV$.
Represent $v$ by $w$ followed by the inclusion into $V$, and
represent $u$ by $w^{-1}\lambda(e)$. Then $uv=1_V$ and
$vu=\lambda(e)$. These operators give a representation of $B$
whose restriction to $S$ is faithful. Thus $B\ne0$ and $S\to B$
is injective. The choice of $w$ proves existence of a model;
it is not needed to compute the finite presentation.
\end{proof}

\begin{corollary}\label{cor:projective-undecidable}
There is no uniform algorithm deciding whether a specified finitely
generated projective over a finitely presented inverse semiring is
zero, even when the additive-idempotent support is $\B$ and the
projective is given by a one-by-one idempotent matrix. The corresponding
isomorphism problem is also undecidable.
\end{corollary}
\begin{proof}
Use Bokut's effective construction
\cite[Section~3, p.~139; Section~4, pp.~142--143]{Bokut} over
$k=\mathbb F_2$. His proof starts with a fixed finitely presented
semigroup $\Pi$ with undecidable word problem and the nonunital free
product $k[\Pi]*ke$, where $e^2=e$ and $ke$ is the chosen
obstruction algebra. Both factors embed by the free-product
description in his Section~2. A difference $a$ of two semigroup
words is zero exactly when the words agree; otherwise it is a
nonzero vector, hence primitive over $k$. His lemma effectively
constructs a finitely presented algebra $A(a)$ that is zero when
$a=0$ and contains the original free product when $a\ne0$.
The named idempotent $e$ therefore vanishes in the first case and
survives in the second. Lemma~\ref{lem:unital-idempotent-bridge}
converts this into a finite unital $\mathbb F_2$-algebra presentation
that is zero exactly when the words agree. This is also a finite
unital ring presentation, with the additional relation $2\cdot1=0$.
Proposition~\ref{prop:projective-decision-reduction} completes the
effective reduction to the claimed projective zero and isomorphism
tests.
\end{proof}
This obstruction concerns finite presentations, not supplied finite
operation tables, and makes no assertion about commutative rings.
It does not make mathematical classification by ordinary-ring data
impossible; it rules out the stated uniform decision procedures.
